% !TeX root = outline.tex
\newlist{paragraphplan}{enumerate}{1}
\setlist[paragraphplan]{label=第\arabic*段,align=left,labelwidth=3em,labelsep=0.6em,leftmargin=3.6em,topsep=0.35em,itemsep=0.3em,parsep=0pt}
\chapter{绪论}
\noindent\textbf{分段大纲}\quad 按章节列出写作内容。物理定义、器件响应关系及尚未落实的研究内容留空。

\section{研究背景与意义}
\subsection{水下目标探测的需求与多物理场途径}
\begin{paragraphplan}
\item 压缩保留海洋开发、安全与水下装备发展所带来的观测需求。
\item 交代水声观测的发展与环境影响，引出磁、电、压力、尾流等观测方向。
\item 按物理来源说明不同特征的观测意义，自然收束到水下电学研究。
\end{paragraphplan}
\subsection{水下电学源的构成与探测困难}
\begin{paragraphplan}
\item 介绍腐蚀、阴极保护、轴系调制、设备杂散和流动感应等电学来源。
\item 从理论计算到实际测量，说明源、介质条件及接收响应之间的关系。
\item 说明空间分布与衰减、位置方向依赖、背景变化及接收漂移带来的问题。
\end{paragraphplan}
\subsection{从测场到测电荷的观测量转变}
\begin{paragraphplan}
\item 区分电位差测量、接收电荷响应及差分电压读出。
\item 明确源端电荷量、空间分布和时间变化是研究对象，介质电位与电场是中间描述量。
\item 区分源端反应与输运、介质空间分布和接收器响应，缺少的物理定义与标定关系留空。
\end{paragraphplan}
\subsection{研究意义}
\begin{paragraphplan}
\item 总述电荷响应研究在水下电学观测中的位置，说明源形成、介质分布与敏感单元输出之间的联系为何值得研究。
\item 展开理论意义。说明电荷源状态、介质场量和接收响应的对应关系如何支持观测解释，深化对电学特征可观测性的认识。
\item 展开实验方法意义。说明等效源如何使源特征可控、实验可重复，并为模型检验、接收标定和不同记录之间的比较提供依据。
\item 展开信号解释意义。说明物理模型与对照数据如何共同支持响应归因，使特征选择具有物理依据，并联系重复性和稳定性评价。
\item 展开感知技术与应用意义。由阵列的时间变化和空间差异过渡到连续状态判别，说明其对海洋观测、水下装备感知及目标探测研究的价值。
\end{paragraphplan}

\section{国内外研究现状}
\subsection{水下电学源机理与建模}
\begin{paragraphplan}
\item 沿海洋腐蚀、电位计算和船体模型交代发展过程。
\item 介绍轴系、设备和流动相关电学源的研究。
\item 比较以介质电位为未知量与以反应输运解释源形成的研究思路，未核实的成果留空。
\item 评述已有源变量与电荷响应观测之间仍需衔接的内容。
\end{paragraphplan}
\subsection{水下电场传感与测量技术}
\begin{paragraphplan}
\item 介绍电极材料研究，分别讨论噪声、漂移、频率响应和灵敏度。
\item 介绍 MEMS 感应式器件，明确本系统技术出处与标定关系的空缺。
\item 介绍多通道及阵列观测的发展，联系空间信息获取与通道一致性。
\item 综合评述不同测量量及其性能比较所需的条件。
\end{paragraphplan}
\subsection{电荷源表征、反演与等效构造}
\begin{paragraphplan}
\item 区分由源预测观测和由观测解释源的问题。
\item 梳理源反演、正则化和等效源表征，尚无依据的内容留空。
\item 介绍受控源实验，再讨论源特征、介质分布及接收输出三个层次的等效对象。
\end{paragraphplan}
\subsection{电学信号处理与特征提取}
\begin{paragraphplan}
\item 保留 FFT、功率谱、小波和模态分解的发展脉络。
\item 介绍电学记录中的去噪、谱线表征与稀疏表示研究。
\item 评述各方法对应的信号信息，并联系电荷响应的实际形态。
\end{paragraphplan}
\subsection{水下目标检测与识别}
\begin{paragraphplan}
\item 介绍统计检测与恒虚警的评价思路。
\item 介绍主动电感知中的样品区分和属性估计，被动实测类别识别内容留空。
\item 梳理变化点与在线判别，尚无依据的内容留空。
\item 区分存在性判断、属性估计与连续状态判别的任务。
\end{paragraphplan}
\subsection{研究现状总结与存在问题}
\begin{paragraphplan}
\item 归纳源状态量与接收观测量的对应问题。
\item 归纳等效构造与接收标定的衔接问题。
\item 归纳设备、水流及背景共同作用下的响应归因问题。
\item 归纳连续判别与参考状态评价的衔接问题。
\end{paragraphplan}

\section{关键科学问题与研究内容}
\subsection{关键科学问题}
\begin{paragraphplan}
\item 源的典型电荷结构及其时变特征如何表征和复现。
\item 源、介质与接收电荷响应如何联系，实际物理关系待明确。
\item 源相关变化如何检出并获得归因支持。介绍有无、方位、距离的任务层级，方位与距离的本文范围留空。
\end{paragraphplan}
\subsection{研究内容}
\begin{paragraphplan}
\item 电荷源表征与等效源构造。
\item 源、介质及接收响应之间的理论联系。
\item 水槽观测、同步记录和接收标定。
\item 特征提取、耦合归因与连续判别，模型与数据在共同观测条件下对应。
\end{paragraphplan}
\subsection{技术路线}
\begin{paragraphplan}
\item 按源表征、观测映射、等效实验、特征分析和判别评价说明研究关系。
\end{paragraphplan}

\section{论文结构安排}

\section{本章小结}
\begin{paragraphplan}
\item 回扣源、映射和判别三个问题，引向典型电荷源的表征与等效。
\end{paragraphplan}
